\section{Construcción del margen previo y definiciones alternativas}
\label{anexo:moderador}

El contraste de heterogeneidad del cuerpo divide a las incumbentes según el
margen con que operaban antes de la entrada. Este anexo detalla cómo se
construye esa medida y por qué se descartan dos definiciones más simples.

\subsection{Por qué una posición relativa y no el margen en niveles}

El objeto que interesa es el poder de mercado local de la estación, que la
\Cref{prop:margen} identifica con el margen de equilibrio $\mu_i = c\,\bar{d}_i$.
El margen en niveles, sin embargo, no es comparable entre estaciones: depende
del combustible, del momento y del mercado en que la estación opera.

La medida empleada es por eso la posición relativa de la estación dentro de la
distribución de precios de su propia comuna en cada mes. Dos propiedades la
hacen apropiada. La primera es que dentro de una comuna y un mes todas las
estaciones enfrentan el mismo costo mayorista, dado que el MEPCO es una
referencia nacional; como el margen porcentual $(p - m)/p$ es estrictamente
creciente en $p$ cuando $p > m > 0$, el orden de los precios coincide
exactamente con el orden de los márgenes. La posición relativa mide entonces el
margen sin necesidad de calcularlo. La segunda es que, al ser un percentil, es
adimensional, lo que permite promediarla entre combustibles cuyos niveles de
precio difieren.

Hay además una razón teórica para preferir la comparación \emph{dentro} de la
comuna. Por la \Cref{prop:margen}, el margen ordena estaciones por aislamiento
espacial solo si el costo de desplazamiento $c$ es común entre ellas. Ese
supuesto es insostenible entre Santiago y una ruta interurbana, pero razonable
dentro de una misma comuna. Comparar posiciones dentro de comuna, y no niveles
entre comunas, es lo que hace que la medida siga leyéndose como aislamiento.

Se exige que la comuna tenga al menos cuatro estaciones en el mes para que el
percentil signifique algo. La restricción tiene un costo: descarta el $23\%$ de
las estaciones de la muestra, con mayor incidencia entre los controles, de modo
que la muestra de este ejercicio es menor que la de la especificación principal.

\subsection{Ponderación entre combustibles}

Una estación vende cuatro combustibles y su poder de mercado no se expresa por
igual en todos. La medida agrega las cuatro posiciones con ponderadores iguales
a la participación de cada combustible en el volumen nacional de ventas que
reporta \textcite{FNE2025informe} para 2017, renormalizados de modo que sumen
uno: $19{,}6\%$ la gasolina de 93 octanos, $11{,}9\%$ la de 95, $4{,}0\%$ la de
97 y $64{,}4\%$ el diésel.

Ponderar por volumen y no por igual evita que la gasolina de 97 octanos, que es
apenas el $4\%$ del mercado, pese lo mismo que el diésel, que es casi dos
tercios. La contrapartida es que los ponderadores provienen de ventas
nacionales, que incluyen canales distintos del minorista —el diésel tiene usos
industriales y de generación eléctrica que no pasan por estaciones de
servicio—, de modo que probablemente sobrerrepresentan el diésel respecto de lo
que vende una estación típica. No existe una desagregación pública por canal que
permita corregirlo.

\subsection{Por qué se excluye el combustible analizado}

El precio observado de una estación puede descomponerse en un componente
persistente, que refleja su aislamiento y es el que interesa, y uno transitorio.
Al seleccionar el tercil superior se seleccionan estaciones que son caras por
ambas razones a la vez. Las que lo son por una perturbación transitoria
favorable revierten después, con entrada o sin ella, de modo que el grupo alto
exhibiría una caída posterior aunque el tratamiento no existiera. Es la forma
que adopta aquí la reversión a la media.

Dos rasgos del diseño la contienen. El primero es que la posición se promedia
sobre toda la ventana previa —cuarenta y ocho meses en la mediana— y no se
toma de un mes aislado, lo que reduce el peso del componente transitorio. El
segundo, y el decisivo, es que el moderador excluye el combustible que se está
analizando: al estimar el efecto sobre el precio del diésel, el moderador se
construye únicamente con las tres gasolinas. La perturbación transitoria que
afectó al diésel de esa estación en la ventana previa no participa entonces de
la clasificación.

Conviene ser preciso sobre el alcance de esta corrección. Elimina el canal en
que la clasificación y el resultado comparten literalmente la misma variable,
pero no aquel en que una perturbación común a la estación mueve todos sus
combustibles a la vez. La correlación entre las posiciones de la gasolina de 93
octanos y del diésel es de $0{,}59$ a nivel de estación y mes, de modo que queda
contaminación residual. La verificación que cierra el argumento es empírica y no
de diseño: las trayectorias previas de ambos grupos son planas, con valores $p$
entre $0{,}18$ y $0{,}92$, lo que es incompatible con una reversión que ya
estuviera en curso antes de la entrada.

\subsection{Comparación con definiciones alternativas}

La \Cref{tab:moderador} reporta la diferencia entre grupos bajo tres
definiciones del moderador, estimadas sobre la misma especificación: el rango del
propio combustible, que es la versión más simple; el compuesto ponderado que
incluye los cuatro combustibles; y el compuesto que excluye el analizado, que es
el utilizado.

\begin{table}[H]
\centering
\caption{Diferencia entre grupos según la definición del margen previo}
\label{tab:moderador}
\tabla{tab_het_moderador}
\end{table}

La comparación entrega tres lecciones.

Primero, agregar combustibles mejora la medida. El paso de la definición A a la
C aumenta la magnitud del contraste y reduce su error estándar en la gasolina de
95 octanos y en el diésel. Promediar cuatro posiciones atenúa el error de
medición del moderador, y una variable de partición menos ruidosa separa mejor a
los grupos.

Segundo, la brecha entre B y C cuantifica el canal mecánico. En el diésel, cuyo
ponderador propio es de $0{,}64$, la definición B arroja una diferencia de
$-0{,}686$ y la C una de $-0{,}562$: cerca de un quinto del contraste bajo B
proviene de incluir el propio combustible en el moderador. Adoptar B produciría
un resultado más nítido a costa de reintroducir justamente el sesgo que se busca
evitar.

Tercero, el caso de la gasolina de 97 octanos es ilustrativo. Bajo la definición
A, que usa su propio rango, el contraste es de $-0{,}357$ y significativo al
$5\%$; bajo la definición C cae a $-0{,}220$ y deja de serlo. La 97 es el
combustible de menor volumen y precio más volátil, es decir, aquel en que el
componente transitorio pesa más, y es por lo tanto donde la definición ingenua
produce el resultado más engañoso. Es la razón concreta por la que este trabajo
no la utiliza.
